% GENERATED FILE. Edit canonical lessons, then run npm run generate:books.
% canonical-source-hash: fnv1a64:576675ae0e260062
% canonical-lessons: FR-C471-deception, FR-C471-decevoir, FR-C471-chagrin, FR-C471-sangloter, FR-C471-consoler

\chapter{Une grande déception}
\label{ch:fr-une-grande-deception}
\hlchaptermodality{471}

\begin{chapteropening}
\textbf{By the end of this chapter:} \emph{I can talk about disappointment and sadness.}

Complete the last lesson of chapter 471: I can talk about disappointment and sadness.
\end{chapteropening}

\section[la déception]{la déception — disappointment}

\label{lesson:FR-C471-deception}



\begin{quote}
Before the new one: say the French for relief, then the French for to reassure.
\end{quote}

\subsection*{You'll want to know: la déception}
\textbf{la déception} — \textquotedblleft{}disappointment\textquotedblright{}.

\begin{quote}
Le film a été une vraie déception.
\end{quote}

\begin{quote}
\emph{The film was a real disappointment.}
\end{quote}

A false friend: it does not mean \emph{deception} in English, which is \emph{la tromperie}.

\begin{cousinweb}
\emph{la déception} is a learned borrowing from Late Latin \emph{dēceptiō} \textquotedblleft{}deceit\textquotedblright{}, from \emph{dēcipere} \textquotedblleft{}to catch, to deceive\textquotedblright{}, built on \emph{capere} \textquotedblleft{}to take\textquotedblright{} (Lewis \& Short).
\end{cousinweb}

\begin{grammarlens}[title={What you've built}]
One more word for telling, explaining and describing at B1.
\end{grammarlens}

\subsection*{Your turn}
\begin{itemize}
\raggedright
  \item \emph{Say it:} \emph{la déception}
  \item \emph{Say it:} \emph{la déception}, once more
  \item \emph{Say it:} a sentence of your own with \emph{la déception}
  \item \emph{Recall:} say the French for relief, then the French for to reassure, then say \emph{la déception} again
\end{itemize}

\begin{tcolorbox}[breakable,colback=teal!4,colframe=teal!35!black,title={Before you move on}]
What does la déception mean? (\textquotedblleft{}Disappointment\textquotedblright{}.) Say it once more.
\end{tcolorbox}
\section[décevoir]{décevoir — to disappoint}

\label{lesson:FR-C471-decevoir}



\begin{quote}
Before the new one: say the French for to reassure, then the French for disappointment.
\end{quote}

\subsection*{You'll want to know: décevoir}
\textbf{décevoir} — \textquotedblleft{}to disappoint\textquotedblright{}.

\begin{quote}
Le concert m'a beaucoup déçu.
\end{quote}

\begin{quote}
\emph{The concert disappointed me a lot.}
\end{quote}

Like \emph{recevoir}: \emph{je déçois}, \emph{j'ai été déçu}. It does not mean to deceive.

\begin{cousinweb}
\emph{décevoir} is inherited, through Old French \emph{deceveir}, from Latin \emph{dēcipere} \textquotedblleft{}to deceive\textquotedblright{} (EtymDB), built on \emph{capere} \textquotedblleft{}to take\textquotedblright{} (Lewis \& Short).
\end{cousinweb}

\begin{grammarlens}[title={What you've built}]
One more word for telling, explaining and describing at B1.
\end{grammarlens}

\subsection*{Your turn}
\begin{itemize}
\raggedright
  \item \emph{Say it:} \emph{décevoir}
  \item \emph{Say it:} \emph{décevoir}, once more
  \item \emph{Say it:} a sentence of your own with \emph{décevoir}
  \item \emph{Recall:} say the French for to reassure, then the French for disappointment, then say \emph{décevoir} again
\end{itemize}

\begin{tcolorbox}[breakable,colback=teal!4,colframe=teal!35!black,title={Before you move on}]
What does décevoir mean? (\textquotedblleft{}To disappoint\textquotedblright{}.) Say it once more.
\end{tcolorbox}
\section[le chagrin]{le chagrin — sorrow, grief}

\label{lesson:FR-C471-chagrin}



\begin{quote}
Before the new one: say the French for disappointment, then the French for to disappoint.
\end{quote}

\subsection*{You'll want to know: le chagrin}
\textbf{le chagrin} — \textquotedblleft{}sorrow, grief\textquotedblright{}.

\begin{quote}
Elle a eu beaucoup de chagrin quand son chat est mort.
\end{quote}

\begin{quote}
\emph{She felt great sorrow when her cat died.}
\end{quote}

\emph{Avoir du chagrin} — to be sad. \emph{Un chagrin d'amour} — heartbreak.

\begin{cousinweb}
The origin of \emph{le chagrin} \textquotedblleft{}sorrow\textquotedblright{} is not settled. A link with \emph{chagrin}, a grainy leather named from Turkish \emph{sağrı} (EtymDB), is sometimes proposed.
\end{cousinweb}

\begin{grammarlens}[title={What you've built}]
One more word for telling, explaining and describing at B1.
\end{grammarlens}

\subsection*{Your turn}
\begin{itemize}
\raggedright
  \item \emph{Say it:} \emph{le chagrin}
  \item \emph{Say it:} \emph{le chagrin}, once more
  \item \emph{Say it:} a sentence of your own with \emph{le chagrin}
  \item \emph{Recall:} say the French for disappointment, then the French for to disappoint, then say \emph{le chagrin} again
\end{itemize}

\begin{tcolorbox}[breakable,colback=teal!4,colframe=teal!35!black,title={Before you move on}]
What does le chagrin mean? (\textquotedblleft{}Sorrow, grief\textquotedblright{}.) Say it once more.
\end{tcolorbox}
\section[sangloter]{sangloter — to sob}

\label{lesson:FR-C471-sangloter}



\begin{quote}
Before the new one: say the French for to disappoint, then the French for sorrow.
\end{quote}

\subsection*{You'll want to know: sangloter}
\textbf{sangloter} — \textquotedblleft{}to sob\textquotedblright{}.

\begin{quote}
L'enfant sanglotait dans les bras de sa mère.
\end{quote}

\begin{quote}
\emph{The child was sobbing in his mother's arms.}
\end{quote}

\emph{Un sanglot} — a sob.

\begin{cousinweb}
\emph{sangloter} is inherited from Latin \emph{singultāre} \textquotedblleft{}to hiccup, to sob\textquotedblright{} (EtymDB), from \emph{singultus} \textquotedblleft{}a hiccup, a sob\textquotedblright{} (Lewis \& Short).
\end{cousinweb}

\begin{grammarlens}[title={What you've built}]
One more word for telling, explaining and describing at B1.
\end{grammarlens}

\subsection*{Your turn}
\begin{itemize}
\raggedright
  \item \emph{Say it:} \emph{sangloter}
  \item \emph{Say it:} \emph{sangloter}, once more
  \item \emph{Say it:} a sentence of your own with \emph{sangloter}
  \item \emph{Recall:} say the French for to disappoint, then the French for sorrow, then say \emph{sangloter} again
\end{itemize}

\begin{tcolorbox}[breakable,colback=teal!4,colframe=teal!35!black,title={Before you move on}]
What does sangloter mean? (\textquotedblleft{}To sob\textquotedblright{}.) Say it once more.
\end{tcolorbox}
\section[consoler]{consoler — to comfort, to console}

\label{lesson:FR-C471-consoler}



\begin{quote}
Before the new one: say the French for sorrow, then the French for to sob.
\end{quote}

\subsection*{You'll want to know: consoler}
\textbf{consoler} — \textquotedblleft{}to comfort, to console\textquotedblright{}.

\begin{quote}
Mon ami m'a consolé après mon échec.
\end{quote}

\begin{quote}
\emph{My friend comforted me after my failure.}
\end{quote}

\emph{Se consoler} — to get over something sad.

\begin{cousinweb}
\emph{consoler} is a learned borrowing from Latin \emph{cōnsōlārī} \textquotedblleft{}to comfort\textquotedblright{} (EtymDB), built on \emph{sōlārī} \textquotedblleft{}to soothe\textquotedblright{} (Lewis \& Short).
\end{cousinweb}

\begin{grammarlens}[title={What you've built}]
One more word for telling, explaining and describing at B1.
\end{grammarlens}

\subsection*{Your turn}
\begin{itemize}
\raggedright
  \item \emph{Say it:} \emph{consoler}
  \item \emph{Say it:} \emph{consoler}, once more
  \item \emph{Say it:} a sentence of your own with \emph{consoler}
  \item \emph{Recall:} say the French for sorrow, then the French for to sob, then say \emph{consoler} again
\end{itemize}

\begin{tcolorbox}[breakable,colback=teal!4,colframe=teal!35!black,title={Before you move on}]
What does consoler mean? (\textquotedblleft{}To comfort, to console\textquotedblright{}.) Say it once more.
\end{tcolorbox}
